\documentclass[a4paper,12pt]{article}
\setcounter{secnumdepth}{5}
\setcounter{tocdepth}{3}
\newcounter{ZhRenew}
\setcounter{ZhRenew}{1}
\newcounter{SectionLanguage}
\setcounter{SectionLanguage}{1}
\input{/usr/share/LaTeX-ToolKit/template.tex}
\title{探索與書寫報告文稿：我對真理的懷疑與重建}
\author{沈威宇}
\date{\temtoday}
\pagestyle{empty}
\begin{document}
\temdoc
\section*{探索與書寫報告文稿：我對真理的懷疑與重建}
\hfill 沈威宇 \temtoday
\subsection{緣起：樸素真理觀的崩解}
高二歷史課，老師在教導歷史之餘總是介紹各家哲學思想，並提出許多發人深省的命題。一天，老師問道「什麼是真理？」我首先直覺地想：「真理就是與事實相符合的命題。太陽從東邊升起、水在攝氏零度結冰，這些在經驗世界中可被反覆驗證的現象，便是真理。」然而這種想法很快就被我自己反駁：「數學對象是否存在？正義又是否具有真實意義？」那一刻，我開始懷疑我是否一直都錯了。
\subsection{衝突：經驗論與形上學的拉扯}
之後的日子裡，我像被兩匹馬往反方向拉扯。一邊是物理實驗的數據、餐廳便當的價格、手機螢幕的文字，這些看得見摸得著的東西，體現出經驗世界的真實。但在另一側，當我擁抱拓樸學那抽象而美麗的數學工具、崇尚《正義論》那平等而公平的美好嚮往，我深知，我終究無法完全拒斥柏拉圖理念世界裡那個完美的三角形。於是，闃寂無聲的班上每頁躊躇，月黑風高的夜裡每步蹣跚，都撕扯著我尚未成形的思想。
\subsection{辯證：實證論的嘗試與挫折}
我不願繼續迷惘，試圖找到一個能夠說服自己的主張。當我讀到艾耶爾的「一個陳述只有在能夠透過分析或經驗驗證的情況下才被認為是字面上的意思。」我感到既興奮又踏實。我開始拿自己的生活做實驗：紀錄天氣、睡眠、讀書、心情等等，企圖用數據分析證明一個陳述要麼是邏輯分析的重言，要麼是物質世界的結果，否則便毫無意義。然而一次在擁擠的捷運上，當我一次次修改 Python 輸出散布圖與迴歸線，我發現變數選擇和降噪算法要比數據本身對結論影響更大，我萌生疑問：怎麼區分觀察術語和理論術語？一個理論要能解釋多少經驗現實才是真理？觀察的選擇、理論的建構、語言的使用，要怎麼擺脫觀察者、語言、環境等造成的相對性？我放下手機，盯著捷運車窗倒映的自己。那張臉皺眉地問我：如果連驗證原則都無法自證，我還能相信什麼？
\subsection{轉折：對數學邏輯的質疑與探索}
轉折的契機發生於我學習數學時。看著書頁上蒼白的文氏圖，我想出羅素悖論，它拉著我一腳躍進了 ZF 集合論的世界。我手舞足蹈，以為終於找到了一套永遠不會出錯的理論、一種可以描述一切的語言。然而當我讀到連續統假設，我又卡住了。它在這個演繹系統內無法被證實或證偽，它獨立於這個理論，卻並非沒有意義——至少，我無法忽視它。夕陽斜照在數學課本上，字跡被染得血紅，我趴在桌上，嘆出的氣吹動書頁一角，連空氣都在訕笑我的徒勞。
\subsection{重建：模型論與證明論框架}
我意識，問題不在於選擇哪一種觀點，而是我需要一種方式，讓不同的「真」可以共存。我開始研讀形式語言、模型論與證明論等知識，漸漸地，我的真理觀被重塑，最終在哲學課的啟發下發展出一個自己的真理框架。我放寬邏輯的選擇、區分經驗與解釋、包容符應的不全，同時接納了經驗論與形上學、容納了多元性與相對性，建立了一個形式化的框架。終於，我不再試圖消弭這樣的張力與不確定性，而是學會承受它、擁抱它。
\subsection{反思：探索仍未結束}
當然，我知道這個框架也充滿漏洞。首先，它太寬容了，幾乎什麼都能被放進來，使真理流於瑣碎。其次，它不對演繹系統作出限制，不保證可操作性。捫心自問，我深知我的探索仍長路漫漫，或許永遠沒有終點，正如沒有一個解釋一切的完美真理一般。現在，如果問我「你相信真理嗎？」我會反問：「你願意和我一起畫一張地圖嗎？」我想，這比告訴他地圖只有一張要誠實得多。
\end{document}
